\chapter{2014年重庆卷（理）}

\section{单选题}

\begin{problem}
在复平面内表示复数 \(\i(1 - 2\i)\) 的点位于（\quad）

\choices
    {第一象限}
    {第二象限}
    {第三象限}
    {第四象限}
\end{problem}

\begin{answer}
A
\end{answer}

\begin{problem}
对任意等比数列 \(\{a_{n}\}\)，下列说法一定正确的是（\quad）

\choices
    {\(a_{1}, a_{3}, a_{9}\) 成等比数列}
    {\(a_{2}\)，\(a_{3}\)，\(a_{6}\) 成等比数列}
    {\(a_{2}\)，\(a_{4}\)，\(a_{8}\) 成等比数列}
    {\(a_{3}\)，\(a_{6}\)，\(a_{9}\) 成等比数列}
\end{problem}

\begin{answer}
D
\end{answer}

\begin{problem}
已知变量 \(x\) 与 \(y\) 正相关，且由观测数据算得样本平均数 \(\bar{x} = 3\)，\(\bar{y} = 3.5\)，则由该观测数据算得的线性回归方程可能是（\quad）

\choices
    {\(\hat{y} = 0.4x + 2.3\)}
    {\(\hat{y} = 2x - 2.4\)}
    {\(\hat{y} = -2x + 9.5\)}
    {\(\hat{y} = -0.3x + 4.4\)}
\end{problem}

\begin{answer}
A
\end{answer}

\begin{problem}
已知向量 \(\bs{a} = (k, 3)\)，\(\bs{b} = (1, 4)\)，\(\bs{c} = (2, 1)\)，且 \((2\bs{a} - 3\bs{b}) \perp \bs{c}\)，则实数 \(k =\)（\quad）

\choices
    {\(-\frac{9}{2}\)}
    {0}
    {3}
    {\(\frac{15}{2}\)}
\end{problem}

\begin{answer}
C
\end{answer}

\begin{problem}
执行如图所示的程序框图，若输出 \(k\) 的值为 6 ，则判断框内可填入的条件是（\quad）

\bitmapfigure[max width=0.42\linewidth]{img/2014/chongqing_science/q5_fig1.png}

\choices
    {\(s > \frac{1}{2}\)}
    {\(s > \frac{3}{5}\)}
    {\(s > \frac{7}{10}\)}
    {\(s > \frac{4}{5}\)}
\end{problem}

\begin{answer}
C
\end{answer}

\begin{problem}
已知命题 \(p\)：对任意 \(x \in \R\)，总有 \(2^x > 0\)； \(q\)： \(x > 1\) 是“\(x > 2\)”的充分不必要条件. 则下列命题为真命题的是（\quad）

\choices
    {\(p \wedge q\)}
    {\(\neg p \land \neg q\)}
    {\(\neg p \land q\)}
    {\(p \wedge \neg q\)}
\end{problem}

\begin{answer}
D
\end{answer}

\begin{problem}
某几何体的三视图如图所示，则该几何体的表面积为（\quad）

\bitmapfigure[max width=0.72\linewidth]{img/2014/chongqing_science/q7_fig1.png}

\choices
    {\(54\)}
    {\(60\)}
    {\(66\)}
    {\(72\)}
\end{problem}

\begin{answer}
B
\end{answer}

\begin{problem}
设 \(F_{1}\)，\(F_{2}\) 分别为双曲线 \(\frac{x^{2}}{a^{2}} - \frac{y^{2}}{b^{2}} = 1\)（\(a > 0, b > 0\)）的左，右焦点，双曲线上存在一点 \(P\) 使得 \(\left|PF_{1}\right| + \left|PF_{2}\right| = 3b\)，\(\left|PF_{1}\right| \cdot \left|PF_{2}\right| = \frac{9}{4} ab\)，则该双曲线的离心率为（\quad）

\choices
    {\( \frac{4}{3} \)}
    {\( \frac{5}{3} \)}
    {\( \frac{9}{4} \)}
    {3}
\end{problem}

\begin{answer}
B
\end{answer}

\begin{problem}
某次联欢会要安排 3 个歌舞类节目，2 个小品类节目和 1 个相声类节目的演出顺序，则同类节目不相邻的排法种数是（\quad）

\choices
    {\(72\)}
    {\(120\)}
    {\(144\)}
    {\(168\)}
\end{problem}

\begin{answer}
B
\end{answer}

\begin{problem}
已知 \(\triangle ABC\) 的内角 \(A\)，\(B\)，\(C\) 满足 \(\sin 2A + \sin (A - B + C) = \sin (C - A - B) + \frac{1}{2}\)，面积 \(S\) 满足 \(1 \leqslant S \leqslant 2\)，记 \(a\)，\(b\)，\(c\) 分别为 \(A\)，\(B\)，\(C\) 所对的边，则下列不等式一定成立的是（\quad）

\choices
    {\(bc(b + c) > 8\)}
    {\(ab(a + b) > 16\sqrt{2}\)}
    {\(6 \leqslant abc \leqslant 12\)}
    {\(12 \leqslant abc \leqslant 24\)}
\end{problem}

\begin{answer}
A
\end{answer}

\section{填空题}

\begin{problem}
设全集 \(U = \{n \in \N \mid 1 \leqslant n \leqslant 10\}\)，\(A = \{1, 2, 3, 5, 8\}\)，\(B = \{1, 3, 5, 7, 9\}\). 则 \((\complement_U A) \cap B =\fillinblank{}\).
\end{problem}

\begin{answer}
\(\{7, 9\}\)
\end{answer}

\begin{problem}
函数 \( f(x) = \log_{2}\sqrt{x} \cdot \log_{\sqrt{2}}(2x) \) 的最小值为\fillinblank{}.
\end{problem}

\begin{answer}
\(-\dfrac{1}{4}\)
\end{answer}

\begin{problem}
已知直线 \(ax + y - 2 = 0\) 与圆心为 \(C\) 的圆 \((x - 1)^2 + (y - a)^2 = 4\) 相交于 \(A\)，\(B\) 两点，且 \(\triangle ABC\) 为等边三角形，则实数 \(a = \)\fillinblank{}.
\end{problem}

\begin{answer}
\(4 \pm \sqrt{15}\)
\end{answer}

\begin{problem}
过圆外一点 \(P\) 作圆的切线 \(PA\)（\(A\) 为切点），再作割线 \(PBC\) 依次交圆于 \(B\)，\(C\)，若 \(PA = 6\)，\(AC = 8\)，\(BC = 9\)，则 \(AB =\fillinblank{}\).
\end{problem}

\begin{answer}
\(4\)
\end{answer}

\begin{problem}
已知直线 \(l\) 的参数方程为 \( \begin{cases}
x=2+t,\\
y=3+t,
\end{cases} \)（\(t\) 为参数），以坐标原点为极点，\(x\) 轴的正半轴为极轴建立极坐标系，曲线 \(C\) 的极坐标方程为 \( \rho\sin^{2}\theta-4\cos\theta=0 \)（\(\rho\geqslant0\)，\(0\leqslant\theta<2\pi\)），则直线 \(l\) 与曲线 \(C\) 的公共点的极径 \( \rho= \)\fillinblank{}.
\end{problem}

\begin{answer}
\(\sqrt{5}\)
\end{answer}

\begin{problem}
若不等式 \( \left|2x-1\right|+\left|x+2\right|\geqslant a^{2}+\frac{1}{2}a+2 \) 对任意实数 \(x\) 恒成立，则实数 \(a\) 的取值范围是\fillinblank{}.
\end{problem}

\begin{answer}
\(\left[-1, \dfrac{1}{2}\right]\)
\end{answer}

\section{解答题}

\begin{problem}
已知函数 \( f(x) = \sqrt{3}\sin (\omega x + \varphi) \)（\( \omega > 0, -\frac{\pi}{2} \leqslant \varphi < \frac{\pi}{2} \)）的图象关于直线 \( x = \frac{\pi}{3} \) 对称，且图象上相邻两个最高点的距离为 \( \pi \).

\begin{enumerate}
\item 求 \( \omega \) 和 \( \varphi \) 的值；
\item 若 \( f\left(\frac{\alpha}{2}\right) = \frac{\sqrt{3}}{4} \)（\(\frac{\pi}{6} < \alpha < \frac{2\pi}{3}\)），求 \( \cos \left(\alpha + \frac{3\pi}{2}\right) \) 的值.
\end{enumerate}
\end{problem}

\begin{solution}
\begin{enumerate}
\item 相邻两个最高点的横坐标相差一个周期，因此
\[
T=\frac{2\pi}{\omega}=\pi,\qquad \omega=2.
\]
图象关于直线 \(x=\dfrac{\pi}{3}\) 对称，故该直线经过图象的最高点或最低点，即
\[
2\cdot\frac{\pi}{3}+\varphi=\frac{\pi}{2}+k\pi\quad\left(k\in\Z\right).
\]
由 \(-\dfrac{\pi}{2}\leqslant\varphi<\dfrac{\pi}{2}\)，得 \(k=0\)，从而 \(\varphi=-\dfrac{\pi}{6}\).

\item 由 \(f\left(\dfrac{\alpha}{2}\right)=\dfrac{\sqrt{3}}{4}\)，得
\[
\sqrt{3}\sin\left(\alpha-\frac{\pi}{6}\right)=\frac{\sqrt{3}}{4},
\qquad \sin\left(\alpha-\frac{\pi}{6}\right)=\frac{1}{4}.
\]
因为 \(\dfrac{\pi}{6}<\alpha<\dfrac{2\pi}{3}\)，所以 \(0<\alpha-\dfrac{\pi}{6}<\dfrac{\pi}{2}\)，从而
\[
\cos\left(\alpha-\frac{\pi}{6}\right)=\frac{\sqrt{15}}{4}.
\]
又 \(\cos\left(u+\dfrac{3\pi}{2}\right)=\sin u\)，所以
\[
\cos\left(\alpha+\frac{3\pi}{2}\right)
 =\sin\alpha
 =\sin\left(\alpha-\frac{\pi}{6}\right)\cos\frac{\pi}{6}
  +\cos\left(\alpha-\frac{\pi}{6}\right)\sin\frac{\pi}{6}
 =\frac{\sqrt{3}+\sqrt{15}}{8}.
\]
\end{enumerate}
\end{solution}

\begin{problem}
一盒中装有9张各写有一个数字的卡片，其中4张卡片上的数字是1，3张卡片上的数字是2，2张卡片上的数字是3，从盒中任取3张卡片.

\begin{enumerate}
\item 求所取 3 张卡片上的数字完全相同的概率；
\item \(X\) 表示所取 3 张卡片上的数字的中位数，求 \(X\) 的分布列与数学期望. （注：若三个数 \(a\)，\(b\)，\(c\) 满足 \(a \leqslant b \leqslant c\)，则称 \(b\) 为这三个数的中位数）
\end{enumerate}
\end{problem}

\begin{solution}
\begin{enumerate}
\item 从 9 张卡片中任取 3 张，共有 \(\mathrm C_9^3=84\) 种取法. 3 张卡片上的数字完全相同的取法数为
\(\mathrm C_4^3+\mathrm C_3^3+\mathrm C_2^3=4+1+0=5\)，所以所求概率为
\[
P=\frac{5}{84}.
\]

\item 当 \(X=1\) 时，取到的数字组合为 \((1,1,1)\)、\((1,1,2)\)、\((1,1,3)\)，对应的取法数为
\[
\mathrm C_4^3+\mathrm C_4^2\mathrm C_3^1+\mathrm C_4^2\mathrm C_2^1
 =4+18+12=34.
\]
当 \(X=3\) 时，取到的数字组合为 \((1,3,3)\)、\((2,3,3)\)，对应的取法数为
\[
\mathrm C_4^1\mathrm C_2^2+\mathrm C_3^1\mathrm C_2^2=4+3=7.
\]
因此 \(X=2\) 的取法数为 \(84-34-7=43\)，故 \(X\) 的分布列为

\begin{center}
\begin{tabular}{c|ccc}
\(X\) & 1 & 2 & 3 \\
\hline
\(P\) & \(\dfrac{17}{42}\) & \(\dfrac{43}{84}\) & \(\dfrac{1}{12}\) \\
\end{tabular}
\end{center}

所以
\[
E(X)=1\cdot\frac{34}{84}+2\cdot\frac{43}{84}+3\cdot\frac{7}{84}
 =\frac{47}{28}.
\]
\end{enumerate}
\end{solution}

\begin{problem}
如图，四棱锥 \(P - ABCD\) 中，底面是以 \(O\) 为中心的菱形，\(PO \perp\) 底面 \(ABCD\)，\(AB = 2\)，\(\angle BAD = \frac{\pi}{3}\)，\(M\) 为 \(BC\) 上一点，且 \(BM = \frac{1}{2}\)，\(MP \perp AP\).

\begin{enumerate}
\item 求 \(PO\) 的长；
\item 求二面角 \(A - PM - C\) 的正弦值.
\end{enumerate}

\bitmapfigure[max width=0.58\linewidth]{img/2014/chongqing_science/q19_fig1.png}
\end{problem}

\begin{solution}
建立空间直角坐标系，使底面 \(ABCD\) 在 \(z=0\) 平面内，并取
\(A=(0,0,0)\)，\(B=(2,0,0)\)，\(D=(1,\sqrt{3},0)\)，则
\(C=(3,\sqrt{3},0)\)，\(O=\left(\dfrac{3}{2},\dfrac{\sqrt{3}}{2},0\right)\).
由 \(M\in BC\)，且 \(BM=\dfrac{1}{2}=\dfrac{1}{4}BC\)，得
\[
M=B+\frac{1}{4}\overrightarrow{BC}
 =\left(\frac{9}{4},\frac{\sqrt{3}}{4},0\right).
\]

\begin{enumerate}
\item 设 \(P=\left(\dfrac{3}{2},\dfrac{\sqrt{3}}{2},h\right)\)，则 \(PO=h\). 由 \(MP\perp AP\)，有
\[
\overrightarrow{MP}=\left(-\frac{3}{4},\frac{\sqrt{3}}{4},h\right),\qquad
\overrightarrow{AP}=\left(\frac{3}{2},\frac{\sqrt{3}}{2},h\right),
\]
\[
\overrightarrow{MP}\cdot\overrightarrow{AP}
 =-\frac{9}{8}+\frac{3}{8}+h^2=0.
\]
所以 \(h^2=\dfrac{3}{4}\)，从而
\[
PO=\frac{\sqrt{3}}{2}.
\]

\item 取平面 \(APM\) 和平面 \(CPM\) 的法向量分别为
\[
\bs{n}_1=\overrightarrow{MP}\times\overrightarrow{MA}
 =\left(\frac{3}{8},-\frac{9\sqrt{3}}{8},\frac{3\sqrt{3}}{4}\right),
\]
\[
\bs{n}_2=\overrightarrow{MP}\times\overrightarrow{MC}
 =\left(-\frac{9}{8},\frac{3\sqrt{3}}{8},-\frac{3\sqrt{3}}{4}\right).
\]
计算得
\[
|\bs{n}_1|^2=\frac{45}{8},\qquad |\bs{n}_2|^2=\frac{27}{8},
\qquad \bs{n}_1\cdot\bs{n}_2=-\frac{27}{8}.
\]
设二面角 \(A-PM-C\) 为 \(\theta\)，则其正弦值与两法向量夹角的正弦值相同，故
\[
\cos\theta=\frac{|\bs{n}_1\cdot\bs{n}_2|}{|\bs{n}_1||\bs{n}_2|}
 =\frac{\sqrt{15}}{5},
\qquad
\sin\theta=\sqrt{1-\frac{15}{25}}=\frac{\sqrt{10}}{5}.
\]
\end{enumerate}
\end{solution}

\begin{problem}
已知函数 \( f(x) = a\e^{2x} - b\e^{-2x} - cx \) （\(a,b,c \in \R\)） 的导函数 \( f'(x) \) 为偶函数，且曲线 \( y = f(x) \) 在点 \(\left(0, f(0)\right)\) 处的切线的斜率为 \(4-c\).

\begin{enumerate}
\item 确定 \(a\)， \(b\) 的值；
\item 若 \(c = 3\)，判断 \( f(x) \) 的单调性；
\item 若 \( f(x) \) 有极值，求 \(c\) 的取值范围.
\end{enumerate}
\end{problem}

\begin{solution}
\begin{enumerate}
\item
\[
f'(x)=2a\e^{2x}+2b\e^{-2x}-c.
\]
由 \(f'(x)\) 为偶函数，得
\[
f'(x)-f'(-x)=2(a-b)\left(\e^{2x}-\e^{-2x}\right)=0.
\]
对任意 \(x\) 恒成立，只能有 \(a=b\). 又切线斜率为
\[
f'(0)=2a+2b-c=4a-c=4-c,
\]
所以 \(a=b=1\).

\item 当 \(c=3\) 时，\(f'(x)=2\e^{2x}+2\e^{-2x}-3\). 由
\(\e^{2x}+\e^{-2x}\geqslant2\)，得 \(f'(x)\geqslant1>0\)，所以 \(f(x)\) 在 \(\R\) 上单调递增.

\item 此时
\[
f'(x)=2\left(\e^{2x}+\e^{-2x}\right)-c=4\cosh 2x-c.
\]
若 \(c\leqslant4\)，则 \(f'(x)\geqslant4-c\geqslant0\)，且当 \(c=4\) 时仅在 \(x=0\) 处取零，导数不发生由正到负或由负到正的改变，故 \(f(x)\) 无极值.
若 \(c>4\)，方程 \(f'(x)=0\) 等价于 \(\cosh 2x=\dfrac{c}{4}>1\)，有两个实根；此时 \(f'(x)\) 在两根外为正、两根之间为负，故 \(f(x)\) 有极大值和极小值.
因此 \(c\) 的取值范围为 \(c>4\).
\end{enumerate}
\end{solution}

\begin{problem}
设椭圆 \(\frac{x^2}{a^2} + \frac{y^2}{b^2} = 1 \) （\(a > b > 0\)）的左，右焦点分别为 \(F_1\)，\(F_2\)，点 \(D\) 在椭圆上，\(DF_1 \perp F_1F_2\)，\(\frac{|F_1F_2|}{|DF_1|} = 2\sqrt{2}\)，\(\triangle DF_1F_2\) 的面积为 \(\frac{\sqrt{2}}{2}\).

\begin{enumerate}
\item 求该椭圆的标准方程；
\item 设圆心在 \(y\) 轴上的圆与椭圆在 \(x\) 轴的上方有两个交点，且圆在这两个交点处的两条切线相互垂直并分别过不同的焦点，求圆的半径.
\end{enumerate}

\bitmapfigure[max width=0.56\linewidth]{img/2014/chongqing_science/q21_fig1.png}
\end{problem}

\begin{solution}
\begin{enumerate}
\item 设椭圆半焦距为 \(c\)，则 \(F_1=(-c,0)\)，\(F_2=(c,0)\)，且
\(c^2=a^2-b^2\). 由 \(DF_1\perp F_1F_2\)，可设 \(D=(-c,y_0)\)，于是
\[
|F_1F_2|=2c,\qquad |DF_1|=|y_0|.
\]
由已知比值得 \(|y_0|=\dfrac{c}{\sqrt{2}}\). 又
\[
\frac{1}{2}|F_1F_2|\cdot|y_0|=\frac{\sqrt{2}}{2},
\]
所以 \(\dfrac{c^2}{\sqrt{2}}=\dfrac{\sqrt{2}}{2}\)，即 \(c^2=1\). 将 \(D\) 代入椭圆方程，得
\[
\frac{c^2}{a^2}+\frac{y_0^2}{b^2}=1,\qquad a^2-b^2=1,\qquad y_0^2=\frac{1}{2}.
\]
令 \(b^2=u>0\)，则 \(a^2=u+1\)，从而
\[
\frac{1}{u+1}+\frac{1}{2u}=1,\qquad 2u^2-u-1=0.
\]
因为 \(u>0\)，得 \(u=1\)，故 \(a^2=2\)，椭圆的标准方程为
\[
\frac{x^2}{2}+y^2=1.
\]

\item 设圆心为 \(Q=(0,k)\)，圆与椭圆在 \(x\) 轴上方的两个交点关于 \(y\) 轴对称，记为 \(T_+=(u,v)\)，\(T_-=(-u,v)\)，其中 \(u>0\)，\(v>0\). 由于两条切线分别过不同焦点，交换两交点的记号后，可设 \(T_+\) 处的切线过 \(F_2=(1,0)\)，\(T_-\) 处的切线过 \(F_1=(-1,0)\).
\[
\overrightarrow{QT_+}=(u,v-k),\qquad \overrightarrow{F_2T_+}=(1-u,-v).
\]
切线与半径垂直，故
\[
u(1-u)-v(v-k)=0,\qquad kv=u^2+v^2-u.
\]
两条切线互相垂直，等价于两条对应半径互相垂直，于是
\[
\overrightarrow{QT_+}\cdot\overrightarrow{QT_-}
 =-u^2+(v-k)^2=0,\qquad k-v=\pm u.
\]
若 \(k=v+u\)，由切线条件得 \(v=u-1\)；若 \(k=v-u\)，由切线条件得 \(v=1-u\). 两种情形都代入椭圆方程
\[
\frac{u^2}{2}+v^2=1.
\]
第一种情形给出 \(u=\dfrac{4}{3}\)，\(v=\dfrac{1}{3}\)；第二种情形给出 \(u=\dfrac{4}{3}\)，\(v=-\dfrac{1}{3}\)，与 \(v>0\) 矛盾. 因此
\[
Q=\left(0,v+u\right)=\left(0,\frac{5}{3}\right).
\]
圆的半径为
\[
r=|QT_+|=\sqrt{u^2+(v-k)^2}
 =\sqrt{\left(\frac{4}{3}\right)^2+\left(-\frac{4}{3}\right)^2}
 =\frac{4\sqrt{2}}{3}.
\]
\end{enumerate}
\end{solution}

\begin{problem}
设 \(a_1 = 1\)，\(a_{n+1} = \sqrt{a_n^2 - 2a_n + 2} + b\)（\(n \in \N^*\)）.

\begin{enumerate}
\item 若 \(b = 1\)，求 \(a_2\)，\(a_3\) 及数列 \(\{a_n\}\) 的通项公式；
\item 若 \(b = -1\)，问：是否存在实数 \(c\)，使得 \(a_{2n} < c < a_{2n+1}\) 对所有 \(n \in \N^*\) 成立？ 证明你的结论.
\end{enumerate}
\end{problem}

\begin{solution}
\begin{enumerate}
\item 当 \(b=1\) 时，
\[
a_2=\sqrt{1^2-2\cdot1+2}+1=2,
\qquad a_3=\sqrt{2^2-2\cdot2+2}+1=1+\sqrt{2}.
\]
下面用数学归纳法证明 \(a_n=1+\sqrt{n-1}\). 当 \(n=1\) 时结论成立；若 \(a_n=1+\sqrt{n-1}\)，则
\[
a_{n+1}=\sqrt{a_n^2-2a_n+2}+1
 =\sqrt{(a_n-1)^2+1}+1
 =\sqrt{n}+1=1+\sqrt{(n+1)-1}.
\]
故对一切 \(n\in\N^*\)，都有 \(a_n=1+\sqrt{n-1}\).

\item 当 \(b=-1\) 时，记
\[
g(x)=\sqrt{(x-1)^2+1}-1,\qquad a_{n+1}=g(a_n).
\]
先算得 \(a_2=g(1)=0<\dfrac{1}{4}\). 对 \(0\leqslant x\leqslant1\)，有
\[
0\leqslant g(x)\leqslant\sqrt{2}-1<1,
\]
所以 \(a_n\in[0,1]\)（\(n\geqslant2\)），且 \(g(x)\) 在 \([0,1]\) 上严格递减. 另外
\[
g\left(\frac{1}{4}\right)=\sqrt{\left(-\frac{3}{4}\right)^2+1}-1=\frac{1}{4}.
\]
因此，若 \(a_{2n}<\dfrac{1}{4}\)，则由 \(g\) 的严格递减性，
\(a_{2n+1}=g(a_{2n})>g\left(\dfrac{1}{4}\right)=\dfrac{1}{4}\)；再由同样的理由，
\(a_{2n+2}=g(a_{2n+1})<\dfrac{1}{4}\). 结合 \(a_2<\dfrac{1}{4}\) 归纳可得
\[
a_{2n}<\frac{1}{4}<a_{2n+1}\qquad\left(n\in\N^*\right).
\]
所以存在这样的实数，取 \(c=\dfrac{1}{4}\).
\end{enumerate}
\end{solution}
